\chapter{Discussion and Limitations}\label{ch:6}

\section{Efficiency gains and scope}\label{sec:6.1}

Both evaluations show lower input use under the compression policy. In the fresh cohort, accuracy also improved over the heuristic at nearly the same input length, which indicates that shorter prompts alone do not explain the difference between those two construction policies. The 300-target comparison establishes the input saving but leaves its accuracy effects unresolved \hyperref[sec:A.2]{[E3, E6]}.

Output-format counts clarify the fresh accuracy result. FULL produced a format-compliant final output on all 246 targets, while PNS did so on 245 and answered 16 more questions correctly than FULL. The heuristic had three format failures, compared with one for PNS. Fewer format failures therefore cannot explain the gain over FULL or the full improvement over the heuristic. The comparison evaluates the complete construction-and-fallback policy, in which reasoning content, candidate search and selection changed together.

The two cohorts address different evaluation settings. The 300-target cohort was baseline-exposed, contained 271 full-CPT groups and included 59 targets sharing full CPTs with prior material. Its clustered analysis accounts for dependence within full-CPT groups; historical exposure remains part of the evaluation setting. The 246 fresh cases test new CPT values within existing graph families and stories under prespecified validity filters. The cohorts are analysed separately, and a difference in statistical significance does not establish a difference in effects. Both evaluations used a fixed 56-identity bank and two same-type demonstrations selected using dataset query-type metadata. Neither evaluation tested performance when query types must be inferred or when the demonstration-routing policy changes.

The fresh evaluation uses the probability-query part of the bank. Its eligible demonstration pool contains the 18 accepted traces of those query types and the applicable parent fallbacks. The 34 accepted deterministic-counterfactual traces contribute to the 52-trace construction result but are outside this evaluation pool.

\section{Semantic validity, recovery and selection}\label{sec:6.2}

Accepted trace q4284 retained the general claim that a causal arrow implies a positive effect, although a causal edge does not specify the sign of an effect \hyperref[sec:A.2]{[E5]}. This example shows how an incorrect causal statement can survive selection when the final label is correct.

The rejected q29833-s029-delete-r4 continuation illustrates a different failure \hyperref[sec:A.2]{[E3]}. Its binary SCM satisfies S = K, M = \ensuremath{\neg}S and Z = S \ensuremath{\land} M. The question asks whether Z = 0 after do(K = 0), which gives S = 0, M = 1 and Z = 0. The 2,073-token candidate answered yes but retained M = 0 and later treated the outcome rule as disjunctive. The original parent already contained the fixed-M premise. The shorter continuation inherited an error rather than showing that DELETE caused it. Label, length and derivation validity therefore require separate records.

A continuation can reconstruct deleted facts or find an alternative solution. A successful DELETE probe therefore shows recovery from an edited context; it does not establish that the deleted step was dispensable in the original derivation. Construction selects parent chains with correct final answers and retains candidates that meet the answer, length and judging criteria. Performance on the construction questions is therefore part of selection, not independent evidence of transfer. The downstream evaluations test whether the selected demonstrations remain useful on other questions.

\section{Construction cost and open questions}\label{sec:6.3}

Reusing the compressed demonstration bank reduced aggregate input use, while total spending also includes bank construction and judging. Completion tokens fell by 0.072\% against FULL in the 300-target cohort and by 11.05\% in the fresh cohort, so completion savings were not proportional to input savings. Phase56 construction recorded 41,392,323 Qwen tokens, alongside additional judging costs and missing usage \hyperref[sec:A.2]{[E3]}. The CNY 28.89 bill covered the complete current two-GPU power-on cycle, including preparation and evaluation, but excluded historical construction, judge charges and unmetered assistance \hyperref[sec:A.2]{[E9]}. Arm-specific charges and an incremental monetary break-even point cannot be calculated from these accounts. Detailed historical and recovery ledgers remain in the electronic appendix.

The immediate comparison to add is budget-matched repeated sampling. KEEP-only or ordinary resampling should match the PNS candidate budget, answer checks, judge access and selection rule. The length-matched heuristic control measures performance against inexpensive shortening but does not separate intervention search from repeated generation and selection. A blinded independent semantic reference covering accepted, rejected and boundary candidates would complement that comparison by assessing validation quality. A further experiment could use the saved input budget for additional demonstrations and compare accuracy at a fixed total prompt length.
